\documentclass[12pt,a4paper]{article}
\usepackage[utf8]{inputenc}
\usepackage[T1]{fontenc}
\usepackage{amsmath}
\usepackage{amssymb}
\usepackage{amsfonts}
\usepackage{mathtools}
\usepackage{physics}
\usepackage{siunitx}
\usepackage{graphicx}
\usepackage{xcolor}
\usepackage{hyperref}
\usepackage{geometry}
\usepackage{microtype}
\usepackage{lmodern}
\geometry{margin=2.5cm}
\hypersetup{colorlinks=true,linkcolor=blue,urlcolor=blue}

\title{td1.pdf}
\author{Jack Ma}
\date{07 mai 2026}

\begin{document}
\maketitle

\section*{td1.pdf}
\section*{TD 1: Espaces de Hilbert et applications}

\section*{Généralités.}

Exercice 1 (Identité du parallélogramme généralisée). Soit ( \(H,\langle\),\(\rangle ) un espace de Hilbert. Soit n \in[|2,+\infty|[\). Soit \(\left(x_{1}, \ldots x_{n}\right) \in H^{n}\). Montrer que l'on a
\[
\sum_{i=1}^{n}\left\|x_{i}\right\|^{2}=\frac{1}{2^{n}} \sum_{\left(\varepsilon_{1}, \ldots, \varepsilon_{n}\right) \in\{-1,1\}^{n}}\left\|\varepsilon_{1} x_{1}+\ldots+\varepsilon_{n} x_{n}\right\|^{2}
\]

Exercice 2. On considère \(\mathbb{R}[X]\) muni du produit scalaire \(\langle P, Q\rangle=\int_{0}^{1} P(x) Q(x) d x\).
1. Démontrer que c'est un produit scalaire.
2. Trouver une suite de polynôme \(\left(P_{n}\right)_{n \in \mathbb{N}}\) qui converge uniformément vers exp sur [ 0,1 ].
3. En déduire \(\left(P_{n}\right)_{n \in \mathbb{N}} \rightarrow \exp\) pour la norme associée au produit scalaire \(\langle\),\(\rangle .\)
4. Que peut-on en déduire sur \((\mathbb{R}[X],\langle\rangle\),\() ?\)

Exercice 3 (Théorème de Fréchet-von Neumann-Jordan). Soit \(E\) un espace de Banach réel muni d'une norme \(\|\cdot\|\) vérifiant l'égalité du parallélogramme:
\[
\|x+y\|^{2}+\|x-y\|^{2}=2\left(\|x\|^{2}+\|y\|^{2}\right), \forall(x, y) \in E^{2} .
\]

On pose
\[
\langle x, y\rangle:=\frac{1}{2}\left(\|x+y\|^{2}-\|x\|^{2}-\|y\|^{2}\right), \forall(x, y) \in E^{2}
\]

On se propose de vérifier que cette expression est un produit scalaire vérifiant de plus \(\langle x, x\rangle=\|x\|^{2}\).
1. Vérifier que \(\langle x, y\rangle=\langle y, x\rangle,\langle-x, y\rangle=-\langle x, y\rangle\), et \(\langle x, 2 y\rangle=2\langle x, y\rangle, \forall(x, y) \in E^{2}\).
2. Montrer que \(\langle x+y, z\rangle=\langle x, z\rangle+\langle y, z\rangle, \forall(x, y, z) \in E^{3}\).
3. Montrer que \(\langle\lambda x, y\rangle=\lambda\langle x, y\rangle, \forall(\lambda, x, y) \in\left(\mathbb{R} \times E^{2}\right)\). On traitera d'abord le cas \(\lambda \in \mathbb{N}\) puis \(\lambda \in \mathbb{Z}\) puis \(\lambda \in \mathbb{Q}\).
4. Conclure.

Exercice 4 (Complexification d'un espace de Hilbert réel). Soit ( \(H,\langle\),\(\rangle ) un espace de Hilbert réel. On considère\) l'espace vectoriel produit \(H^{\mathbb{C}}:=H \times H\), où l'on pose par définition, pour \((x, y) \in H \times H, i .(x, y):=(-y, x)\).
1. Trouver une addition et une loi de multiplication externe qui permettent de munir \(H^{\mathbb{C}}\) d'une structure d'espace vectoriel complexe. Démontrer que la restriction de cette structure aux réels coïncide avec la structure d'espace vectoriel produit sur \(H \times H\).
2. On identifie \(H\) à \(H \times\{0\}\). Identifier \(i H\) à un sous-espace vectoriel réel de \(H^{\mathbb{C}}\). Montrer qu'en tant qu'espace vectoriel réel, \(H^{\mathbb{C}}\) est la somme directe de \(H\) et \(i H\).
3. On considère \(\left(z_{1}, z_{2}\right) \in H^{\mathbb{C}} \times H^{\mathbb{C}}\), que l'on décompose de manière unique sous la forme \(z_{j}=x_{j}+i y_{j}, j=1,2\). On pose alors
\[
\left\langle z_{1}, z_{2}\right\rangle_{H^{\mathrm{c}}}:=\left\langle x_{1}, x_{2}\right\rangle+\left\langle y_{1}, y_{2}\right\rangle+i\left(\left\langle x_{1}, y_{2}\right\rangle-\left\langle x_{2}, y_{1}\right\rangle\right)
\]

Montrer que c'est l'unique produit scalaire hermitien qui prolonge le produit scalaire sur \(H\).
L'espace \(H^{\mathbb{C}}\) est appelé le complexifié de \(H\).

\section*{Théorème de projection sur un convexe fermé.}

Exercice 5. On considère l'espace \(E=\left(C^{0}([-1,1], \mathbb{R}),\|\cdot\|_{\infty}\right)\), et \(F\) le sous-espace vectoriel formé des fonctions impaires dont l'intégrale est nulle sur \([0,1]\). Soit \(\varphi: t \in[-1,1] \mapsto t\).
1. Montrer que \(F\) est fermé.
2. Montrer que \(d(\varphi, F) \geqslant 1 / 2\).
3. Existe-t-il \(\psi \in F\) tel que \(\|\varphi-\psi\|_{\infty}=1 / 2\) ?
4. Montrer que \(d(\varphi, F)=1 / 2\). Indication: essayer "d'approcher" \(t-1 / 2\) par des fonctions de \(F\).

Commenter.
Exercice 6. On considère l'espace \(E=\left(C^{0}([-1,1], \mathbb{R}),\|\cdot\|_{\infty}\right)\) et \(D\) la droite engendrée par \(x \mapsto 1-x\). Montrer que la distance de \(D\) à la fonction \(x \mapsto 1\) est atteinte en plusieurs points. Commenter.

Exercice 7. Soit \((H,\langle\rangle\),\() un espace de Hilbert. Calculer la projection sur la boule unité fermée.\)
Exercice 8. Soit \((H,\langle\rangle\),\() un espace de Hilbert. Démontrer que tout convexe fermé admet un unique élément de\) norme minimale.

Exercice 9. Soit \(H=l^{2}(\mathbb{N}, \mathbb{R})\) muni du produit scalaire canonique. On note \(C:=\left\{\left(x_{n}\right)_{n \in \mathbb{N}} \mid x_{n} \geqslant 0\right\}\). Montrer que \(C\) est un convexe fermé et calculer la projection orthogonale sur \(C\).

Exercice 10 (Convexes emboîtés et projection, I). Soit \(\left(C_{i}\right)_{i \in \mathbb{N}^{*}}\) une suite de convexes fermés d'un espace de Hilbert \(H,\langle\),\(\rangle tels que C_{1} \supset C_{2} \ldots \supset C_{n} \supset \ldots\). On pose \(C_{\infty}=\bigcap_{n=1}^{\infty} C_{n}\) et on suppose que \(C_{\infty} \neq 0\). On note \(P_{i}\) la projection sur le convexe fermé \(C_{i}\).
1. Montrer que \(C_{\infty}\) est un convexe fermé. On note \(P_{\infty}\) le projection sur le convexe \(C_{\infty}\).
2. Pour \(h \in H\), on pose \(a_{i}=P_{i}(h)\). Démontrer que la suite \(\left\|h-a_{i}\right\|_{i \in \mathbb{N}^{*}}\) est croissante et majorée.
3. Montrer que \(a_{i} \rightarrow P_{\infty}(h)\) quand \(i \rightarrow \infty\). On commencera par montrer que \(\left(a_{i}\right)_{\in \mathbb{N}}\) est une suite de Cauchy, en utilisant la question précédente.

Exercice 11 (Convexes emboîtés et projection, II). Soit \(\left(C_{i}\right)_{i \in \mathbb{N}^{*}}\) une suite de convexes fermés d'un espace de Hilbert \(H,\langle\),\(\rangle tels que C_{1} \subset C_{2} \ldots \subset C_{n} \subset \ldots\). On pose \(C_{\infty}=\bigcup_{n=1}^{\infty} C_{n}\). On note \(P_{i}\) la projection sur le convexe fermé \(C_{i}\).
1. Montrer que \(C_{\infty}\) est un convexe fermé. On note \(P_{\infty}\) le projection sur le convexe \(C_{\infty}\).
2. Pour \(h \in H\), on pose \(a_{i}=P_{i}(h)\). Montrer que \(a_{i} \rightarrow P_{\infty}(h)\) quand \(i \rightarrow \infty\). On commencera par montrer que \(\left(a_{i}\right)_{\in \mathbb{N}}\) est une suite de Cauchy.
Exercice 12 (Théorème de Hahn-Banach, version géométrique et Hilbert). On considère \(C\) un convexe fermé non vide d'un espace de Hilbert réel ( \(H,\langle\),\(\rangle ).\)
1. Soit \(x \in H\). Démontrer qu'il existe \(f \in H^{\prime}\) et \(\alpha \in \mathbb{R}\) tels que
\[
f(x)<\alpha<f(y), \forall y \in C
\]

Que ceci signifie-t-il géométriquement?
2. En déduire que tout convexe fermé propre (i.e. non égal à \(H\) tout entier) s'écrit comme une intersection de demi-espaces fermés.
3. Démontrer que tout sous-espace vectoriel fermé strict s'écrit comme intersection d'hyperplans fermés.
4. Soit \(\widehat{C}\) un autre ensemble convexe non vide, supposé compact, tel que \(C \cap \widehat{C}=\emptyset\). Démontrer qu'il existe \(f \in H^{\prime}\) tel que
\[
\sup _{x \in C} f(x)<\inf _{y \in \widehat{C}} f(y)
\]

Que ceci signifie-t-il géométriquement?
Cette propriété reste-t-elle vraie si on suppose \(\widehat{C}\) seulement fermé?

Exercice 13 (Espérance conditionnelle). On considère un espace probabilisé ( \(\Omega, \mathcal{A}, P\) ) et \(\mathcal{G}\) une sous-tribu de \(\mathcal{A}\). Soit \(X \in L^{1}(\Omega, \mathbb{R})\). Démontrer les différentes propriétés suivantes de l'espérance conditionnelle:
1. Si \(Z\) est \(\mathcal{G}\)-mesurable et essentiellement bornée, alors \(\mathbb{E}[Z E[X \mid \mathcal{G}]]=Z \mathbb{E}[X \mid \mathcal{G}]\).
2. \(\mathbb{E}[\mathbb{E}[X \mid \mathcal{G}]]=\mathbb{E}[X]\).
3. Si \(\mathcal{H}\) est une sous-tribu de \(\mathcal{G}\), alors \(\mathbb{E}[\mathbb{E}[X \mid \mathcal{G}] \mid \mathcal{H}]=\mathbb{E}[X \mid \mathcal{H}]\).
4. Si \(Z\) est \(\mathcal{G}\)-mesurable et essentiellement bornée, alors \(\mathbb{E}[X Z \mid \mathcal{G}]=Z \mathbb{E}[X \mid \mathcal{G}]\).
5. Si \(X\) est \(\mathcal{G}\)-mesurable, alors \(\mathbb{E}[X \mid \mathcal{G}]=X\).
6. Si \(\varphi\) est une fonction convexe telle que \(\varphi(X)\) est intégrable alors \(E[\varphi(X) \mid \mathcal{G}] \geqslant \varphi(E[X \mid \mathcal{G}])\).

\section*{Théorème de Riesz.}

Exercice 14. Soit \(E\) l'espace des suites complexes \(\left(u_{n}\right)_{n \in \mathbb{N}^{*}}\) nulles à partir d'un certain rang, muni du produit scalaire hermitien usuel.
1. Montrer que l'application \(\varphi: u=\left(u_{n}\right) \in E \mapsto \sum_{n=1}^{\infty} \frac{u_{n}}{n}\) est une forme linéaire continue sur \(E\).
2. Existe-t-il un élément \(a \in E\) tel que \(\varphi(u)=\langle a, u\rangle\) ? Que peut-on en déduire sur \(E\) ?

Exercice 15 (Théorème de Radon-Nykodym, version faible). Soit ( \(X, \mathcal{T}\) ) un espace mesurable, soit \(\mu\) et \(\nu\) deux mesures finies sur \(X\). On suppose que quel que soit le mesurable \(A\), on a \(\nu(A) \leqslant \mu(A)\).
1. Démontrer que pour toute fonction \(h\) mesurable positive, on a \(\int_{X} h(x) d \nu(x) \leqslant \int_{X} h(x) d \mu(x)\).
2. En déduire que \(L^{2}(\mu) \subset L^{2}(\nu)\) et que pour tout \(g \in L^{2}(\mu)\) on a \(\|g\|_{L^{2}(\nu)} \leqslant\|g\|_{L^{2}(\mu)}\).
3. Montrer que \(\varphi: g \in L^{2}(\mu) \mapsto \int_{X} g(x) d \nu(x)\) est une forme linéaire continue.
4. En déduire qu'il existe une fonction \(f\) mesurable telle que \(\nu\) soit la mesure de densité de \(f\) par rapport à \(\mu\), i.e. pour tout mesurable \(A\), on a \(\nu(A)=\int_{A} f(x) d \mu(x)\).
Exercice 16 (Universalité de la convolution). Soit \(f\) une fonction de \(R^{N}\left(N \in \mathbb{N}^{*}\right)\) à valeurs réelles. Pour \(x \in \mathbb{R}^{N}\), on pose \(\tau_{x} f: y \mapsto f(y-x)\). Un opérateur \(T\) agissant sur des fonctions de \(L^{2}\left(\mathbb{R}^{N}\right)\) est dit invariant par translation si \(T\left(\tau_{x} f\right)=\tau_{x}(T f)\).

On considère dorénavant \(T: L^{2}\left(\mathbb{R}^{N}\right) \rightarrow \mathcal{C}_{b}\left(\mathbb{R}^{N}\right)\) un opérateur linéaire continu et invariant par translation.
1. Montrer que \(f \mapsto T f(0) \in L^{2}\left(\mathbb{R}^{N}\right)^{\prime}\).
2. En écrivant que \(T f(x)=\tau_{-x}(T f)(0)\), en déduire l'existence de \(g \in L^{2}\left(\mathbb{R}^{N}\right)\) tel que \(T f(x)=f * g(x)\).

Exercice 17 (Adjoint d'un opérateur). Soit \(u\) un endomorphisme continu d'un espace de Hilbert réel ou complexe \((H,\langle\rangle\),\() et y \in H\).
1. Montrer qu'il existe un unique \(z \in H\) tel que pour tout \(x \in H\), on ait \(\langle u(x), y\rangle=\langle x, z\rangle\).

On pose \(u^{*}: y \mapsto z\) défini précédemment. \(u^{*}\) est appelé adjoint de \(u\).
2. Démontrer que \(u^{*}\) est un opérateur linéaire continu et que \(\left\|u^{*}\right\| \leqslant\|u\|\).
3. Montrer que \(u^{* *}=u\).
4. Montrer que \(\left\|u^{*}\right\|=\|u\|\).
5. Montrer que l'application \(u \in \mathcal{L}_{c}(H) \mapsto u^{*} \in \mathcal{L}_{c}(H)\) est anti-linéaire, continue, bijective d'inverse continu.
6. Soit \(v\) un autre endomorphisme continu de \(H\). Montrer que \((u \circ v)^{*}=v^{*} \circ u^{*}\).
7. Montrer que \(\left\|u^{*} \circ u\right\|=\|u\|^{2}=\left\|u \circ u^{*}\right\|\).
8. Montrer que \(\operatorname{Ker}\left(\mathrm{u}^{*}\right)=\operatorname{Im}(\mathrm{u})^{\perp}\). Que vaut \(\operatorname{Ker}\left(\mathrm{u}^{*}\right)^{\perp}\) ?

Exercice 18 (Calculs de quelques adjoints). 1. On considère \(H=l^{2}(\mathbb{N}, \mathbb{C})\) et \(\left(a_{n}\right)_{n \in \mathbb{N}}\) une suite bornée de nombres complexes. On considère \(T:\left(u_{n}\right)_{n \in \mathbb{N}} \mapsto\left(a_{n} u_{n}\right)_{n \in \mathbb{N}}\). Montrer que \(T\) est linéaire continu de \(H\) dans \(H\) et calculer son adjoint.
2. On considère \(H=l^{2}(\mathbb{N}, \mathbb{C})\). On considère \(S:\left(u_{n}\right)_{n \in \mathbb{N}}=\left(u_{0}, u_{1}, \ldots\right) \mapsto\left(v_{n}\right)_{n \in \mathbb{N}}:=\left(0, u_{0}, u_{1}, \ldots\right)\). Montrer que \(S\) est linéaire continu de \(H\) dans \(H\) et calculer son adjoint.
3. On considère \(H=L^{2}([0,1], C)\) et \(K:[0,1] \times[0,1] \rightarrow \mathbb{C}\) une fonction continue. On considère \(U: f \in H \mapsto \left(x \mapsto \int_{0}^{1} K(x, y) f(y) d y\right)\). Montrer que \(U\) est linéaire continu de \(H\) dans \(H\) et calculer son adjoint.

Exercice 19 (Triplet de Gelfand canonique). On considère ( \(\Omega, \mathcal{T}, \mu\) ) un espace mesuré. On considère \(H=L_{\mu}^{2}(\Omega, \mathbb{R})\). On se donne \(m\) une fonction mesurable telle qu'il existe \(\delta>0\) tel que pour presque tout \(x \in \Omega\) on a \(m(x) \geqslant \delta\), et telle que \(m(x)<\infty\) presque partout. On considère \(\left\{=f\right.\) mesurable \(\left.\mid \int_{\Omega} m f^{2} d \mu<\infty\right\}\).
1. Vérifier que \(V \subset H\) avec inclusion continue et \(V\) est dense dans \(H\).
2. Identifier \(V^{\prime}\) avec un espace simple.

\section*{Orthogonalité et bases hilbertiennes.}

Exercice 20. Soit \(E\) l'espace des suites complexes \(\left(u_{n}\right)_{n \in \mathbb{N}^{*}}\) nulles à partir d'un certain rang, muni du produit scalaire hermitien usuel. On considère l'application \(\varphi: u=\left(u_{n}\right) \in E \mapsto \sum_{n=1}^{\infty} \frac{u_{n}}{n}\), qui est une forme linéaire continue sur \(E\).
1. Montrer que \(\operatorname{Ker}(\varphi)\) est un hyperplan fermé, et que \(\operatorname{Ker}(\varphi)^{\perp}=\{0\}\).
2. De manière plus générale, si \((H,\langle\rangle\),\() est un espace préhilbertien (réel ou complexe) non complet, démontrer\) qu'il existe un hyperplan fermé dont l'orthogonal est réduit à 0 . Indication: on admettra que \(H\) peut être plongé dans un espace de Hilbert \((\widehat{H},\langle\rangle\),\() , et que H\) est dense dans \(\widehat{H}\).

Exercice 21. On considère \(H=l^{2}(\mathbb{N}, \mathbb{C})\). On fixe \(n \in \mathbb{N}^{*}\) et on considère \(M=\left\{\left(x_{k}\right)_{k \in \mathbb{N}} \in H \mid \sum_{k=0}^{n} x_{k}=0\right\}\).
1. Montrer que \(M\) est un sous-espace vectoriel fermé de \(H\).
2. Trouver le supplémentaire orthogonal de \(M\).
3. Calculer la distance de la suite ( \(1,0,0, \ldots\) ) à \(M\).

Exercice 22 (Caractérisation des projecteurs orthogonaux). Soit un espace de Hilbert réel ( \(H,\langle\),\(\rangle ) et p \in \mathcal{L}_{c}(H)\) un projecteur de \(H\), i.e. un endomorphisme continu vérifiant \(p \circ p=p\).
1. Démontrer que \(H\) s'écrit comme somme directe de \(\operatorname{Im}(p)\) et \(\operatorname{Ker}(p)\).
2. Démontrer que \(p\) est un projecteur orthogonal \(\Leftrightarrow p\) est 1 -Lipschitzien. Indication: pour le sens \(\Leftarrow\), on fera un dessin et regarder un \(x \in(\operatorname{Ker}(p))^{\perp}\).
Exercice 23. Soit \(E=C^{0}([0,1], \mathbb{R})\), muni du produit scalaire \(L^{2}\). Soit \(C=\{f \in E \mid f(0)=0\}\). Montrer que \(F^{\perp}=\{0\}\). Commenter.

Exercice 24. Soit \((H,\langle\rangle\),\() un espace de Hilbert séparable. Montrer que toute famille orthonormale peut se compléter\) en une base hilbertienne de \(H\).

Exercice 25. 1. Calculer
\[
\min _{a, b, c} \int_{-1}^{1}\left|x^{3}-a x^{2}-b x-c\right|^{2} d x
\]
puis trouver
\[
\max \int_{-1}^{1} x^{3} g(x) d x
\]
parmi les \(g \in L^{2}(-1,1)\) soumis à la contrainte
\[
\int_{-1}^{1} g(x) d x=\int_{-1}^{1} x g(x) d x=\int_{-1}^{1} g(x) d x=0 ; \int_{-1}^{1}|g(x)|^{2} d x=1
\]
2. Soit ( \(H,\langle\),\(\rangle ) un espace de Hilbert. Soit x_{0} \in H\) et \(M\) un sous-espace vectoriel fermé de \(H\), énoncer le problème de maximisation associé au problème \(\min _{x \in M}\left\|x_{0}-x\right\|\) sur le modèle de la question précédente.

Exercice 26. Soit \(H=L^{2}(\mathbb{R})\), et \(V=\left\{f \in C_{c}^{\infty}(\mathbb{R}) / \int_{\mathbb{R}} f(t) d t=0\right\}\).
On rappelle que \(C_{c}^{\infty}(\mathbb{R})\) désigne les fonctions \(C^{\infty}\) à support compact (i.e. nulles en dehors d'un compact), que \(L^{2}(\mathbb{R})\) est un espace de Hilbert pour le produit scalaire usuel et que \(C_{c}^{\infty}(\mathbb{R})\) est dense dans \(L^{2}(\mathbb{R})\).
1. Soit \(\phi \in C_{c}^{\infty}(\mathbb{R})\) telle que \(\int_{\mathbb{R}} \phi(t) d t=1\). On pose
\[
\forall(n, t) \in \mathbb{N}^{*} \times \mathbb{R}, \phi_{n}(t)=\frac{1}{n} \phi\left(\frac{t}{n}\right)
\]
(a) Montrer que
\[
\phi_{n} \underset{n \rightarrow+\infty}{\rightarrow} 0 \text { dans } H
\]
(b) Soit \(g \in C_{c}^{\infty}(\mathbb{R})\). On pose
\[
h=g-\left(\int_{\mathbb{R}} g(t) d t\right) \phi_{n}
\]

Montrer que la fonction \(h\) appartient au sous-espace \(V\).
2. En déduire que
\[
V^{\perp} \subset C_{c}^{\infty}(\mathbb{R})^{\perp}
\]
3. Conclure que le sous-espace \(V\) est dense dans \(H\).

Exercice 27 (Théorème d'échantillonnage de Shannon). On considère
\[
B L^{2}=\left\{f \in L^{2}(\mathbb{R}) \mid \widehat{f}=0 \text { sur } \mathbb{R} \backslash[-1 / 2,1 / 2]\right\}
\]
1. Montrer que \(B L^{2}\) est un sous-espace fermé de \(L^{2}\). On considérera donc maintenant que \(B L^{2}\) est un espace de Hilbert pour le produit scalaire \(L^{2}\).
2. Montrer que toute fonction \(f \in B L^{2}\) est continue et de plus \(\|f\|_{\infty} \leqslant\|f\|_{2}\).
3. Calculer la transformée de Fourier de la fonction indicatrice de \([-1 / 2,1 / 2]\). On appellera cette fonction sinc.
4. on pose \(\tau_{x} f: y \mapsto f(y-x)\). Montrer que \(\left\{\tau_{k}(\operatorname{sinc}) \mid \mathrm{k} \in \mathbb{Z}\right\}\) est une base hilbertienne de \(B L^{2}\). Interpréter ce résultat. Que donne la formule de Parseval dans ce cas-là?
5. On décompose \(f \in B L^{2}\) comme \(f(x)=\sum_{k \in \mathbb{Z}} f_{k} \tau_{k}(\operatorname{sinc})(\mathrm{x})\). Démontrer que la convergence de cette série est uniforme.
6. Démontrer que tout \(f \in B L^{2}\) est une fonction entière.

Exercice 28 (Polynômes orthogonaux). Soit \(I\) un intervalle de \(\mathbb{R}\). On appelle fonction poids une fonction \(\rho\) mesurable, strictement positive, telle que \(\int_{I}|x|^{n} \rho(x)<\infty, \forall n \in \mathbb{N}\). On note \(L^{2}(I, \rho)\) l'espace des fonctions de carré intégrale pour la mesure de densité \(\rho\), muni de son produit scalaire canonique
\[
\langle f, g\rangle:=\int_{I} \rho(x) f(x) \bar{g}(x) d x
\]

On rappelle que cet espace est un espace de Hilbert.
1. Montrer que \(L^{2}(I, \rho)\) contient l'ensemble des polynômes.
2. Montrer qu'il existe une unique famille de polynômes unitaires \(\left(P_{n}\right)_{n \in \mathbb{N}}\) orthogonale et telle que pour tout \(n \in \mathbb{N}, \operatorname{deg}\left(P_{n}\right)=n\).
3. Démontrer que les zéros des polynômes \(P_{n}\) sont distincts, réels et tous dans l'intervalle \(I\). On introduira le polynôme \(S(x)=\prod_{i=1}^{m}\left(x-x_{i}\right)\), où \(x_{i}\) sont les points à l'intérieur de \(I\) où \(P_{n}\) change de signe.
4. Démontrer que pour \(n \geqslant 1\) on a la formule par récurrence \(P_{n+1}=\left(x-A_{n}\right) P_{n}-B_{n} P_{n-1}\), où
\[
A_{n}=\frac{\left\langle x P_{n}, P_{n}\right\rangle}{\left\langle P_{n}, P_{n}\right\rangle} \text { et } B_{n}=\frac{\left\langle x P_{n}, P_{n-1}\right\rangle}{\left\langle P_{n-1}, P_{n-1}\right\rangle}
\]
5. Démontrer par récurrence que pour tout \(n \in \mathbb{N}\), on a
\[
P_{n+1}^{\prime}(x) P_{n}(x)-P_{n}^{\prime}(x) P_{n+1}(x)>0
\]

En déduire qu'entre (au sens strict) chaque racine du polynôme \(P_{n}\) se trouve exactement une racine de \(P_{n+1}\).
6. Exemple (polynômes d'Hermite): on pose \(I=\mathbb{R}\) et \(\rho(x)=e^{-x^{2}}\) (c'est bien une fonction poids).
(a) Calculer \(P_{0}, P_{1}, P_{2}, P_{3}\).
(b) Montrer que pour tout \(n \in \mathbb{N}^{*}\), on a
\[
P_{n}(X)=\frac{(-1)^{n}}{2^{n}} e^{x^{2}}\left(\frac{\mathrm{~d}^{\mathrm{n}}}{\mathrm{dx}^{\mathrm{n}}} e^{-x^{2}}\right) .
\]
(c) Montrer que \(\left\|P_{n}\right\|_{2}=\pi^{1 / 4} \sqrt{n!}\).

Soit \(\xi \in \mathbb{R}, n \in \mathbb{N}\) et
\[
\forall x \in \mathbb{R}, P_{n}^{\xi}(x):=\sum_{k=0}^{n} \frac{(-i \xi x)^{k}}{k!}
\]
(d) Montrer que \(\forall x \in \mathbb{R}, P_{n}^{\xi}(x) \underset{n \rightarrow+\infty}{\rightarrow} e^{-i \xi x}\).
(e) Montrer que \(\forall(n, x) \in \mathbb{N} \times \mathbb{R},\left|P_{n}^{\xi}(x)\right| \leq e^{|\xi x|}\).
(f) Soit \(f \in L^{2}(\mathbb{R})\). On suppose que : \(\forall n \in \mathbb{N},\left\langle P_{n}, f\right\rangle=0\). Soit \(g(x)=e^{-\frac{x^{2}}{2}} f(x)\). Montrer que \(g \in L^{1}(\mathbb{R})\).
(g) Montrer que
\[
\forall \xi \in \mathbb{R}, \hat{g}(\xi)=\int_{\mathbb{R}} e^{-i \xi x} g(x) d x=0
\]

Indication. On pourra calculer les intégrales \(\int_{\mathbb{R}} P_{n}^{\xi}(x) g(x) d x\).
(h) En déduire que \(f\) est identiquement nulle.
(i) Donner une base hilbertienne de \(L^{2}\left(\mathbb{R}, e^{-x^{2}}\right)\).
7. Déduire de l'exemple précédent une base Hilbertienne de \(L^{2}(\mathbb{R})\).

\section*{Autres applications.}

Exercice 29. Soit \(H=\ell^{2}(\mathbb{N}, \mathbb{R})\). On rappelle que \(H\) est un espace de Hilbert pour le produit scalaire:
\[
\forall(x, y) \in H^{2},<x, y>_{H}=\sum_{n=0}^{+\infty} x_{n} y_{n}
\]

Les fonctions bilinéaires suivantes sont-elles continues et coercives sur \(H\) ?
(i) \(a(x, y)=\sum_{n=0}^{+\infty} x_{n} y_{n+1}\), (ii) \(b(x, y)=\sum_{n=0}^{+\infty} x_{n+1} y_{n+1}\),
(iii) \(c(x, y)=x_{0} y_{0}+\sum_{n=0}^{+\infty}\left(x_{n+1}-x_{n}\right)\left(y_{n+1}-y_{n}\right)\), (iv) \(d(x, y)=\sum_{n=0}^{+\infty}\left(x_{n+1} y_{n+1}+2 x_{n} y_{n+1}+2 x_{n} y_{n}\right)\).

Exercice 30 (Théorème de Babuska-Lax-Milgram, 1971). On considère deux espaces de Hilbert réels \((U,\langle\rangle\), et \((V,(\mid))\) et \(a: U \times V\) une forme bilinéaire continue, i.e. \(\exists \beta>0\) tel que pour tout \((u, v) \in U \times V\), on a \(|a(u, v)| \leqslant \beta\|u\|_{U}\|v\|_{V}\). On suppose de plus que \(a\) est faiblement coercive au sens suivant: il existe \(\alpha>0\) tel que \(\sup _{\|v\|=1}|a(u, v)| \geqslant c| | u| |\) et pour tout \(v \neq 0 \in H\), on a \(\sup _{\|u\|=1}|a(u, v)|>0\). Soit \(l \in V^{\prime}\).

Démontrer qu'il existe un unique \(u \in U\) tel que \(a(u, v)=l(v), \forall v \in V\).

Exercice 31 (Equation de Schrödinger). Soit \(L>0\). On considère l'équation suivante, appelée équation de Schrödinger avec conditions de Dirichlet homogènes
\[
\left\{\begin{aligned}
i \partial_{t} u+\partial_{x x} u & =0 & & \text { dans }] 0, \\
u(0, x) & =u_{0}(x) & & \in C^{4}[0, L] \\
u(t, x=0) & =u(t, x=L)=0 & & \text { sur }[0, \infty[
\end{aligned}\right.
\]

Une solution de (1) est une fonction \(u \in C^{1}\left([0, T] \times[0, L]\right.\) telle que pour tout \(t \geqslant 0\), on a \(x \mapsto u(t, x) \in C^{2}([0, L])\), vérifiant (1).
1. Chercher les solutions de cette équation à variables séparées.
2. Démontrer l'existence d'une solution.
3. Montrer que toute solution \(u\) à (1) est telle que son énergie est conservée:
\[
\int_{0}^{L} u^{2}(t) d t=\int_{0}^{L} u_{0}^{2}(x) d x
\]

On pourra multiplier la première ligne de (1) par \(\bar{u}\), intégrer sur \([0, L]\) puis effectuer une intégration par parties.
4. En déduire l'unicité des solutions.
5. Démontrer que \(u\) est en fait définie sur \(\mathbb{R} \times[0, L]\). Cette propriété est appelée réversibilité en temps.

Exercice 32 (Equation de la chaleur \(2 D\) ). Soit \(L_{1}>0\) et \(L_{2}>0\). Considère l'équation suivante, appelée équation de la chaleur sur un rectangle avec conditions de Dirichlet homogènes
\[
\left\{\begin{aligned}
\partial_{t} u-\partial_{x x} u-\partial_{y y} u & =0 & & \text { sur }] 0,+\infty[\times] 0, L_{1}[\times] 0, L_{1}[ \\
\lim _{t \rightarrow 0^{+}} u(0, x) & =u_{0}(x) & & \in L^{2}\left(\left[0, L_{1}\right] \times\left[0, L_{2}\right]\right) \\
u(t, 0, y) & =u\left(t, L_{1}, y\right)=0 & & \text { sur }] 0, \infty\left[\times\left[0, L_{2}\right]\right. \\
u(t, x, 0) & =u\left(t, x, L_{2}\right)=0 & & \text { sur }] 0, \infty\left[\times\left[0, L_{2}\right]\right.
\end{aligned}\right.
\]

Sur le modèle de ce qui a été fait en cours pour l'équation de la chaleur, démontrer l'existence d'une unique solution (dans une classe appropriée) pour cette équation. Indication: on démontrera que sur les fonctions à variables \(x\) et \(y\) séparées sont denses dans \(L^{2}(] 0, L_{1}[\times] 0, L_{2}[)\).

Exercice 33 (Base orthonormale d'ondelettes de Haar). On considère
\[
\psi(x):\left\{\begin{array}{l}
x \in[0,1 / 2[\mapsto 1 \\
x \in[-1 / 2,1] \mapsto-1 \\
x \notin[0,1] \mapsto 0
\end{array}\right.
\]

Pour \(j \in \mathbb{Z}\) et \(k \in \mathbb{Z}\), on pose
\[
\psi_{j, k}(x):=\frac{1}{2^{\frac{j}{2}}} \psi\left(\frac{x-2^{j} k}{2^{j}}\right) .
\]

Le but est est de démontrer que les \(\psi_{j, k}\) forment une base Hilbertienne de \(L^{2}(\mathbb{R})\). On note par la suite \(W_{j}\) l'espace des fonctions de \(L^{2}(\mathbb{R})\) constantes sur les intervalles de la forme \(2^{j}+z\) avec \(z \in \mathbb{Z}\) et de moyenne nulle.
1. Montrer que les \(\psi_{j, k}\) forment une famille orthonormale de \(L^{2}(\mathbb{R})\).
2. Donner une base de \(W_{j}\).
3. Démontrer que les fonctions de \(L^{1}(\mathbb{R}) \cap L^{2}(\mathbb{R})\) à moyenne nulle sont denses dans \(L^{2}(\mathbb{R})\).

Indication: pour \(f \in L^{1}(\mathbb{R}) \cap L^{2}(\mathbb{R})\), on pourra s'intéresser pour \(R>0\) à \(f-\frac{\int_{\mathbb{R}} f}{R} \chi_{[0, R]}\).
4. En déduire que les \(\psi_{j, k}\) forment une base hilbertienne de \(L^{2}(\mathbb{R})\).


\end{document}
